\begin{table}[H]
\centering
\caption{Effect of NIS2 on Enterprise Cybersecurity Investment}
\label{tab:main}
\begin{threeparttable}
\begin{tabular}{lcccc}
\toprule
 & (1) & (2) & (3) & (4) \\
 & DiD & DiD & DDD & Dosage \\
\midrule
Medium $\times$ Post & 0.42 & 0.42 & -0.33 & 0.42 \\
 & (0.80) & (0.80) & (0.94) & (0.80) \\
Medium $\times$ Post $\times$ Transposed & & & 3.35^{***} & \\
 & & & (1.19) & \\
Large $\times$ Post & & & & 2.23^{***} \\
 & & & & (0.63) \\
\midrule
Country FE & Yes & & & \\
Size FE & Yes & & & \\
Year FE & Yes & & & \\
Country $\times$ Size FE & & Yes & Yes & Yes \\
Country $\times$ Year FE & & Yes & Yes & Yes \\
Observations & 162 & 162 & 162 & 243 \\
$R^2$ (within) & 0.001 & 0.002 & 0.026 & 0.043 \\
\bottomrule
\end{tabular}
\begin{tablenotes}[flushleft]
\small
\item \textit{Notes:} Dependent variable is the overall cybersecurity security index (average of 15 individual measures, in percentage points). Medium firms (50--249 employees) are newly regulated under NIS2; small firms (10--49) are exempt. Column (3) interacts with an indicator for countries that completed NIS2 transposition by the October 2024 deadline (Belgium, Croatia, Hungary, Italy, Latvia, Lithuania). Column (4) adds large firms ($\geq$250 employees), already regulated under NIS1 and subject to intensified NIS2 requirements. Standard errors clustered at the country level in parentheses. * $p<0.10$, ** $p<0.05$, *** $p<0.01$.
\end{tablenotes}
\end{threeparttable}
\end{table}
